\section{Conclusion}
\label{sec:conclusion}

VRASection bridges the gap between geographic compactness and Voting Rights
Act compliance in congressional redistricting.

The core tension in algorithmic VRA redistricting is that race-neutral
algorithms satisfy \emph{Shaw}'s non-targeting requirement but may produce
VRA-deficient maps; while race-targeting algorithms (MetisVra) may satisfy
Section 2 but fail \emph{Shaw}'s predominance test --- and, empirically,
fail numerically in 60\,\% of tested states.

VRASection resolves this tension through a design insight grounded in
Gingles Prong 1 itself: a minority community that satisfies the legal
standard for Section 2 protection is, by definition, one that is
\emph{geographically compact}.
Geographic compactness is already encoded in the census-tract adjacency
graph as clustering.
VRASection makes this clustering measurable through the alignment score
$A(\text{split}) = |\mathrm{MVAP}_\text{frac}(L) - \mathrm{MVAP}_\text{frac}(R)|$,
which detects natural geographic boundaries between minority-concentrated
and minority-dispersed areas without racial targeting.

The tie-breaking approach --- applying the alignment bonus only at the
first bisection level, bounded to 40\,\% of the selection score --- produces
maps that are legally defensible from both directions:
they satisfy Section 2 by exploiting the same geographic concentration that
Gingles Prong 1 requires, and they satisfy \emph{Shaw} by keeping geographic
compactness as the predominant factor with a quantifiable 20\,pp margin.

The Alabama empirical result demonstrates the approach's practical significance.
VRASection shifts the first-level bisection from $1:6$ (Birmingham isolation)
to $2:5$ (Black Belt concentration), at a 4.3\,\% EC premium.
This shift responds directly to the \emph{Allen v.\ Milligan} (2023) finding:
the Black Belt region, dispersed under the $1:6$ split, is concentrated on
one side under the $2:5$ split, creating the geographic preconditions for
a second Black-opportunity district.

The full six-state empirical comparison, pending completion, will test whether
VRASection achieves the MetisVra target (MM district count) without the
MetisVra failure modes (numerical instability, 0\,\% success in three states).
We expect VRASection to succeed where MetisVra failed because it operates
through geographic signal extraction rather than constraint solving,
and to require no user-specified racial targets because the geographic
concentration signal is self-contained in the minority VAP distribution.

More broadly, VRASection demonstrates that the Gingles framework and the
Shaw framework are not in conflict for algorithmically drawn maps.
Gingles Prong 1's geographic compactness requirement is the same
geographic compactness that \emph{Shaw} protects as a traditional redistricting
principle.
An algorithm that rewards natural geographic concentration --- not racial
identity, but the spatial clustering that makes a community legally cognizable
--- respects both requirements simultaneously.
VRASection provides a mathematically precise, legally articulable, and
empirically testable mechanism for doing so.

\paragraph{Future work.}
Open questions for follow-on research include:
(1) validating the alignment score as a proxy for Gingles Prong 1 by
comparing $A(\text{best split})$ to Moran's I spatial autocorrelation
across the seven-state dataset;
(2) evaluating ``VRA-aware full recursion'' (applying alignment at multiple
recursion levels) and characterising its legal risk;
(3) extending VRASection to multi-member districts for state legislative
chambers;
and (4) integrating VRASection output with the Callais evidence layer
for the full two-stage map-generation-plus-evidence workflow.
